% ===== §28 =====
\section{The Capture of Intimate Culture by Capital}
\label{p24:sec:capture}

\noindent
Production having been laid out, the argument turns to capture: to the means by which capital draws value from culture, and, in the end, from the intimate culture the second part described. The account proceeds through three historical stages, each a deeper reach of capital into culture, and the third is the one that touches intimate relation directly. The stages are not offered as a complete history but as a deepening, each stage taking as its raw material something the previous stage left untouched, until the last reaches the bond itself.

\subsection*{Stage one: the culture industry}

The first stage is the commodification of culture as such, analyzed by Adorno and Horkheimer under the name of the \emph{culture industry}. Culture, in their account, becomes an industry in the strict sense: its products are manufactured on the model of factory production, standardized and streamlined, and the individuality they appear to offer is a \emph{pseudo-individuality}, a spurious uniqueness stamped onto a standardized product to simulate a diversity that the underlying standardization has abolished \parencite{adornohorkheimer2002}. What is captured, in this stage, is the aesthetic longing itself: the desire that, in the terms of this paper, reaches toward the Real, toward the impossible object around which the drive turns, is intercepted and filled with a standardized commodity that offers itself as the object of that longing. In the terms the series has developed, this is a degradation at the level of value: a real value, oriented to the unsymbolizable, is forced into the form of a symbolic-commodity value that can be exchanged, and the longing is drawn off into the purchase. The culture industry captures culture by manufacturing its objects and selling them to a desire it has first standardized.

\subsection*{Stage two: the society of the spectacle}

The second stage, analyzed by Debord as the \emph{society of the spectacle}, is a deeper reach, for it captures not merely cultural objects but social relations themselves. In the spectacle, capital no longer only sells cultural commodities; it organizes the whole of social life into a vast accumulation of spectacles, so that the relation between people is mediated by images, and lived directly relations give way to relations between images of relations \parencite{debord1994}. What is captured here is relation itself, in its represented form: the direct coupling that this paper places at the center of intimate culture is replaced, in the spectacle, by a mediated pseudo-relation, a relation to the image of a bond in place of the bond. This is a decisive step toward the capture of intimate culture, for it installs, between the two who would generate a culture of their own, the mediating image supplied by capital, and it teaches the couple to relate to the representation of their relation---to how their bond appears, to its image displayed and consumed---in place of the relation itself.

\subsection*{Stage three: affective and attentional capitalism}

The third stage is the one that reaches intimate relation directly, and it is the living frontier of the present. Contemporary platform capitalism takes intimacy, companionship, and affective connection themselves as raw material to be extracted. The forms are familiar and multiplying: dating platforms that render the search for a bond into a monetized funnel, engineered to be entered and re-entered rather than concluded; social platforms that turn the display of intimacy into a stream of engagement to be harvested; the quantification of affective life into metrics of participation, likes, and attention, so that the bond itself becomes a source of data and revenue. In the terms of political economy, this is the \emph{real subsumption} of the most private region of human life: the drive's interweaving in the couple, which the second part showed to be the very substance of intimate culture, is here taken up as the raw material of an extraction, drilled into as the richest available source of value precisely because it lies nearest the constitutive lack. This stage is the direct capture of the intimate, and it operates, as the phenomenology's final section showed, through the colonization of the couple's own co-signification: it supplies the scripts, owns the platforms, and shapes the sense of what a bond should be, so that the couple, generating what they take to be their own culture, generate value for the owner of the field.

\subsection*{The mechanism of capture, and an independent treatment}

The precise mechanism by which capital captures relational value---the formal chain by which a relational property is stripped from its relation, reified into an attribute, bound to a commodity, and made to yield extraction through the necessary failure of the object of desire---has been analyzed at length by the present author in an independent work, to which the reader is referred and on which this section does not depend \parencite{wanhongrelationalvalue}. That work stands on its own and is not a part of the present argument; it is cited here because it anatomizes, in detail this section need not repeat, the psychoanalytic structure of the capture whose cultural form is here at issue. The point to be retained for the present argument is the one the psychoanalytic line has insisted on throughout: that the capture reaches its object through the drive. Capital captures intimate culture most powerfully because the interwoven drive that charges that culture lies nearest the constitutive lack, and to divert a desire that circles so near the empty place mobilizes far more than any surface preference; the intimate bond is thus the richest fuel for capture, and the most destroyed by it, since the ore is destroyed in the extracting.

\subsection*{Why the capture cannot be total}

The section closes by marking the limit that the framework of Part One established and that the next sections will develop. The capture of intimate culture, however deep, cannot be total, and the reason is the one the second revision of the framework secured: the Real is not completely symbolizable, and the relational Real at the core of the bond therefore resists the symbolic capture that commodification is. What capital can capture is the symbolic surface of intimate culture---the forms, the scripts, the displays, the data---and what it cannot capture is the relational Real those forms served, because that Real is not a symbolizable content that could be extracted but the unsymbolizable core of the bond itself. The capture, at its most complete, is therefore a \emph{decoherence}: it preserves the sign of the bond and draws off its substance, keeping the outward cultural form while the real coherence that charged it is lost.

Here the paper advances a third thesis of its own, one that distinguishes its account of capture from the account of the capture of relational value given in the author's independent work. Call it \emph{capture-as-decoherence}: capital's capture of a relational culture succeeds not by seizing the culture itself but by decohering it---by preserving the symbolizable symbolic form while drawing off its coherence in the register of the Real, so that what remains is a shell whose charge has been withdrawn. Where the independent treatment analyzed capture as a rerouting of desire toward a substitute object, the mechanism marked here is the extraction of coherence: the sign of the bond is kept in circulation and made to yield value, while the real coherence it once carried is dissolved in the keeping. This is why the capture can approach completeness in the symbolic register and yet meet, in the register of the Real, a limit it cannot pass: the symbolic form, being symbolizable, is available to be preserved and sold, and the real coherence, not being symbolizable, is not available to be extracted but only to be lost. Decoherence names exactly this asymmetry, that the capture takes everything of the culture that can be represented and destroys, rather than acquires, what cannot---so that the market ends holding the form of a bond whose substance its own operation has dissolved.

This is why the captured intimate culture is a shell, and why the capture, for all its power, meets a limit it cannot cross. The couple that has surrendered its culture to the market retains, in the relational Real that no market can reach, the possibility of generating again a culture of its own---a possibility the final sections must now examine, first in the constraint that bounds but does not determine, then in the ethics of a cultural practice that would reclaim what capture draws off.
